\documentclass[11pt]{article}

\usepackage[T1]{fontenc}
\usepackage[utf8]{inputenc}
\usepackage{amsmath,amssymb,amsthm,mathtools}
\usepackage[a4paper,margin=30mm]{geometry}
\usepackage{microtype}
\usepackage{xurl}
\usepackage[hidelinks]{hyperref}

\newtheorem{theorem}{Theorem}[section]
\newtheorem{lemma}[theorem]{Lemma}
\newtheorem{proposition}[theorem]{Proposition}
\newtheorem{definition}[theorem]{Definition}
\newtheorem{remark}[theorem]{Remark}

\newcommand{\N}{\mathbb N}
\newcommand{\WIN}{\mathrm{WIN}}
\newcommand{\LOSS}{\mathrm{LOSS}}
\newcommand{\DRAW}{\mathrm{DRAW}}

\title{Termination under Optimal Play in the Two-Player $3n\mathbin{\pm}1$ Game}
\author{Grisha Pochuev\\
  \small\href{mailto:n_854@mail.ru}{n\_854@mail.ru}\\
  \small\url{https://github.com/Grisha-Pochuev/3n-plus-minus-1-game}}
\date{Preprint, 8 August 2026}

\begin{document}
\maketitle

\begin{center}
\small This preprint is distributed under the Creative Commons Attribution
4.0 International license (CC BY 4.0).
\end{center}

\begin{abstract}
We consider the impartial two-player game on positive odd integers in which a
move replaces $n>1$ by the odd part of either $3n-1$ or $3n+1$, and the player
who reaches $1$ wins. Arbitrary play need not terminate, so the relevant
question is whether a draw can survive optimal play. We give a binary normal
form with one expanding branch and one strictly contracting branch. The global
argument marks finite outcome proofs by their canonical heights, ranks finite
multisets of these tokens by an ordinal, and proves well-foundedness of the
remaining equal-rank routing fibres. An exhaustive entry analysis then turns
any hypothetical draw continuation into an infinite path in a well-founded
marked transition system. Consequently no draw positions exist and optimal
play reaches $1$ from every odd starting value. A detailed proof supplement,
regression tests, and the exact status of the ongoing Lean formalization are
included in the accompanying repository. The repository also contains a
declarative global-routing certificate, a small independent checker, and a
Lean proof of the general well-founded-certificate metatheorem. Their exact
trust boundary is stated explicitly: the global assembly is machine-checked
conditional on four local refinement obligations that remain part of the
human proof.
\end{abstract}

\section{Introduction}

The game studied here appears as Conway's \emph{Beans-Don't-Talk} game in
Guy's accounts \cite{Guy1986,Guy1996} and more recently as the $3n{+}{-}1$
game of Alth\"ofer, Hartisch, and Zipproth
\cite{AlthoferHartischZipproth2024}. Its similarity to Collatz iteration is
superficial in one important respect: the two signs are strategic choices by
alternating players.

The specific problem addressed here is the two-player problem stated on
Alth\"ofer's Collatz-problems page: prove that every odd starting value reaches
$1$ when both players act optimally \cite{AlthoferProblemPage}. The formulation
on that page and the game rules used below agree exactly.

For the player to move, a position is $\WIN$ if some move leads to a $\LOSS$
position, and is $\LOSS$ if every move leads to a $\WIN$ position. Because the
game graph is infinite, these two inductive classes need not exhaust all
positions. We call every remaining position $\DRAW$. Thus finite-game backward
induction cannot be used without a separate well-foundedness argument.

Indeed arbitrary play may cycle; for example $5\to7\to5$ is possible. The
main result concerns optimal play and is therefore strictly stronger than the
existence of some terminating choice and different from termination under all
choices.

\begin{theorem}[Main theorem]\label{thm:main}
For every positive odd starting integer, optimal play reaches $1$ after
finitely many moves. Equivalently, the game has no $\DRAW$ positions.
\end{theorem}

\paragraph{Independent verification package.}
The complete proof supplement and the independently runnable verification
artifacts are maintained at
\begin{center}
\url{https://github.com/Grisha-Pochuev/3n-plus-minus-1-game}.
\end{center}
An external reader can clone the repository and run \texttt{python audit.py}
to check the documented arithmetic regressions, finite proof certificates,
and the conditional global-routing assembly. The same repository contains the
negative tests and Lean-checked certificate metatheory, together with an
explicit account of what remains human-verified.

This manuscript presents the final dependency chain. Complete local case
proofs are contained in the numbered proof supplement
\cite{PochuevRepository}; section numbers below refer to that supplement.
The supplement, executable checks, certificate inventory, and formalization
sources are openly available at
\url{https://github.com/Grisha-Pochuev/3n-plus-minus-1-game}.

\section{Rules and binary normal form}

For a positive integer $x$, write $\operatorname{odd}(x)$ for the result of
dividing $x$ by the largest power of two that divides it. From an odd $n>1$
the two children are
\[
  \operatorname{odd}(3n-1),\qquad \operatorname{odd}(3n+1).
\]
The state $1$ is terminal and losing for the player to move.

Write $n=2m+1$ and define
\[
 F(m)=\left\lceil\frac{3m}{2}\right\rceil.
\]
Let $R(m)$ be the integer obtained by deleting the maximal alternating suffix
of the binary expansion of $m$. For instance,
$R(1001_2)=10_2$ and $R(10101_2)=0$.

\begin{proposition}[First normal form]\label{prop:first-normal}
For $m>0$, the two children in $m$-coordinates are exactly
\[
   F(m),\qquad F(R(m)).
\]
\end{proposition}

\begin{proof}
We have $3n-1=2(3m+1)$ and $3n+1=2(3m+2)$. Exactly one of
$3m+1,3m+2$ is odd. Splitting on the final bit of $m$ gives $F(m)$ from
that expression. In the other expression, each additional division by two
crosses one alternating bit of $m$; the first repeated adjacent bit stops the
division. This is precisely deletion of the maximal alternating suffix. The
four prefix/final-bit cases are given in the supplement's normal-form proof.
\end{proof}

Every child in Proposition~\ref{prop:first-normal} lies in the image of $F$.
Representing $F(q)$ by $q$ gives the conjugated game
\[
 A(q)=F(q)=\left\lceil\frac{3q}{2}\right\rceil,
 \qquad B(q)=R(F(q)),
\]
with terminal state $q=0$.

\begin{proposition}[Contracting branch]\label{prop:B-contracts}
For every $q>0$, $B(q)<q$.
\end{proposition}

\begin{proof}
Deleting a nonempty binary suffix gives
\[
B(q)\le \left\lfloor\frac{F(q)}2\right\rfloor
\le \left\lfloor\frac{3q+1}{4}\right\rfloor<q
\]
for $q\ge2$, while $B(1)=0$ directly.
\end{proof}

The difficulty is that $A$ expands and the length of the suffix deleted by
$R$ is unbounded. In particular no fixed residue modulus determines $R$.

\section{Finite proof tokens}

The ordinary inductive outcome construction gives every certified $\WIN$ or
$\LOSS$ position a finite canonical proof height. Terminal state $0$ has
height zero. A $\WIN$ certificate selects a lower-height $\LOSS$ child; a
$\LOSS$ certificate contains both lower-height $\WIN$ children.

During the draw analysis we retain the exact certified endpoints used by
these finite proofs. A marked configuration therefore carries a finite
multiset $M$ of proof heights. These are \emph{tokens}: their identities and
source provenance are retained, not merely their largest height or their
cardinality.

\begin{lemma}[Token multiset rank; Supplement Section 129]
For a finite multiset $M$ of proof heights define
\[
  \Phi(M)=\mathop{\#}_{h\in M}\omega^h,
\]
where $\#$ denotes natural ordinal sum. Replacing one token of height $h$ by
finitely many certified descendants of heights strictly below $h$ strictly
decreases $\Phi$.
\end{lemma}

The use of a multiset, rather than a scalar maximum, is essential: a routing
step may split one obligation into several lower ones. Equality of $\Phi$
alone is not used to identify configurations. The transition inventory proves
that every token-changing edge is strict, hence an equal-$\Phi$ route
preserves the actual token multiset.

\section{Marked factor routing}

Long constant binary tails admit coordinates
\[
 Q_D^g(C)=C2^D-g,
 \qquad g\in\{0,1\},
\]
with $C$ odd. The minimum-source analysis produces marked factor scans in
these coordinates. Their role is to keep a finite outcome witness attached
while the arithmetic route is normalized.

\begin{lemma}[Factor attachment; Supplement Sections 130--135]
Every marked factor scan retains its incoming finite proof tokens. For a
marked source $b=Q_D^g(C)$ with $D\ge3$, let $y$ be its marked $\LOSS$ child
and $q$ its marked $\WIN$ child. There is a common $\WIN$ child $r$ with
\[
  h(r)\le h(y)-1.
\]
Before an inner branch is followed, the scan replaces the carried occurrence
of $y$ by $r$. Every exit retains this lower token and its exact source
attachment. The remaining rows $D=1,2$ are exact lifts over an already marked
token and create no unranked fresh witness.
\end{lemma}

This statement closes a common logical gap: reaching a later finite endpoint
is insufficient unless it is comparable with the witness already used in the
outer rank. Sections 130--135 prove precisely that comparison.

\section{Well-founded fibres}

We use the following standard projection principle in a form adapted to the
marked relation.

\begin{lemma}[Well-founded fibre lemma; Supplement Section 132]
Let a transition relation carry a projection into a well-founded order.
Suppose the projection never increases, every projection-changing edge is
strictly decreasing, and the transition relation restricted to each fixed
projection fibre is well-founded. Then the full transition relation is
well-founded.
\end{lemma}

The proof is well-founded induction on the projection, followed by induction
inside the current fibre. The formulation permits inner finite heights that
are incomparable between different routes; only the retained outer projection
must be monotone.

\begin{lemma}[Fixed-fibre normalizer; Supplement Section 136]
Fix the retained coefficient source and the outer proof-token multiset. The
complete generic obligation/factor routing relation is well-founded.
\end{lemma}

\begin{proof}[Proof architecture]
Canonical $A$-streaks have a four-side-test exit. Valuation-two $B$ transfers
install and then lower a finite local witness. Higher valuations either lower
the source or enter a factor fork that pays in source or proof height. A high
return re-enters only the marked factor rows of the preceding section.

Within the fibre, the internal source/token measure has the form
\[
  \omega^\omega a+\Psi(N),
\]
where $a$ is the inner source and $N$ its finite token multiset. Every return
from the generic subsystem crosses a strict high-return replacement before a
new marked pair is installed. Consequently cross-transitions cannot alternate
forever between two separately terminating subsystems. Applying the fibre
lemma to the complete transition inventory proves the claim. All arithmetic
subcases are expanded in Section 136 and its cited lemmas.
\end{proof}

\section{Entry completeness}

Two analyses feed the normalizer. Sections 14--27 treat a least coefficient
source through constant-tail frames. Sections 69--90 treat height one
separately; its resonance is coupled across consecutive factor levels and is
not replaced by a fixed-depth computation.

\begin{lemma}[Entry trichotomy; Supplement Section 137]
Every outcome-compatible continuation from the minimum-source and height-one
analyses has one of three outputs:
\begin{enumerate}
  \item strict coefficient-source descent;
  \item strict descent of the retained proof-token multiset rank;
  \item entry into the fixed-fibre normalizer of Section 136.
\end{enumerate}
There is no fourth arithmetic exit.
\end{lemma}

The supplement audits the earlier case splits against this trichotomy. This is
the exhaustiveness step that connects the local arithmetic to the global
well-founded relation.

\section{Exclusion of draws}

\begin{proof}[Proof of Theorem~\ref{thm:main}]
Assume a $\DRAW$ exists and choose one with least coefficient source $s$.
Following the exhaustive outcome diamonds produces a marked macro-configuration
that still contains an actual $\DRAW$. By the entry trichotomy, every
macrostep either reaches a $\DRAW$ below $s$, strictly decreases the retained
token rank $\Phi$, or stays within a fixed normalizer fibre.

The first alternative contradicts minimality of $s$. The second cannot occur
infinitely because ordinals are well-founded. The third cannot occur
infinitely by the fixed-fibre normalizer lemma. The well-founded fibre lemma
therefore shows that the complete marked transition relation admits no
infinite path.

On the other hand, a $\DRAW$ position has no $\LOSS$ child and has at least one
$\DRAW$ child. Repeatedly following such a child through the finite
macro-normalizations would give an infinite marked transition path, a
contradiction. Thus no $\DRAW$ exists.

Every state is consequently $\WIN$ or $\LOSS$ with finite canonical proof
height. Optimal play from a $\WIN$ state selects a lower-height $\LOSS$ child,
and every move from a $\LOSS$ state reaches a lower-height $\WIN$ child. The
height decreases until conjugated state $0$, corresponding to original state
$1$.

Finally, if an original odd state is divisible by three, then after either
legal move the odd part of $3n\pm1$ is not divisible by three and lies in the
conjugated image. Hence the conclusion covers every odd starting state.
\end{proof}

\section{Verification status and reproducibility}

The repository distinguishes four kinds of evidence.

First, the mathematical supplement labels each claim as proved,
computationally verified, conjectural, disproved, or an unverified lead.
Second, standard-library Python code checks the exact normal form, descent
identities, routing arithmetic, and bounded retrograde soundness. These finite
runs are regression evidence and are not used as the proof of the infinite
theorem.

Third, the repository contains a finite declarative inventory of the global
routing assembly. It has 15 control states, 46 transition schemas, and 12
guard partitions, including symbolic unbounded valuation and tail-length
intervals. An independent standard-library Python checker verifies source
integrity, complete coverage of the \emph{declared} cases, one rule per case,
lexicographic reset discipline, and acyclicity of the 17 equal-rank control
edges. The remaining 29 schemas have an explicit strict rank component.

The checker prints the status \texttt{CONDITIONAL\_MACHINE\_CHECK}. Four
obligations deliberately remain on the human side of the trust boundary:
\begin{enumerate}
  \item the universal arithmetic and coordinate identities attached to each
        macro schema;
  \item refinement of the declared semantic guards to every legal game case;
  \item outcome compatibility of the selected DRAW continuation and the
        separately carried finite tokens;
  \item finite productivity of every macro-normalization.
\end{enumerate}
Thus acceptance of the JSON certificate validates the final size-change
assembly conditional on the cited Sections 14--137; it is not presented as a
complete machine proof from the original move relation.

Fourth, a pinned Lean 4 project proves the general metatheory used by the
certificate: the lexicographic product of well-founded relations is
well-founded, a relation mapped to strict decreases in such a rank is
well-founded, and a well-founded certified relation admits no infinite route.
The Lean sources use no problem-specific axioms and contain no
\texttt{sorry} or \texttt{admit}. Lean does not yet prove the concrete
multiset rank or the four certificate-to-game refinement obligations, and the
coverage table states this limitation explicitly.

This separation is part of the result: it permits an auditor to reproduce the
computations without confusing them with the well-foundedness proof, and to
measure future formalization progress without importing the desired theorem
as an axiom.

From a clean checkout, the complete locally machine-checkable audit is:
\begin{verbatim}
python audit.py
\end{verbatim}
It runs the complete unit and regression test suite, checks the documented finite
arithmetic range, validates the global assembly certificate, generates a
finite outcome proof DAG, and validates that DAG with a separate checker. The
global certificate alone is checked by
\begin{verbatim}
python scripts/verify_global_certificate.py
\end{verbatim}
The Lean metatheory is checked independently by \texttt{lake build} in the
\texttt{formal} directory. Exact commands, expected output, negative tests,
and the human audit route are recorded in the repository's
\texttt{START\_HERE.md} and \texttt{AUDIT.md} files.

\section{Known proof hazards}

The argument does not propagate a least draw indefinitely along the expanding
branch, infer a theorem from a search depth, use a fixed modulus for an
unbounded suffix, or infer winning play from a merely terminating one-player
strategy. It also does not replace retained witnesses by incomparable later
witnesses. Most importantly, the fixed-fibre proof checks transitions between
the obligation and factor subsystems; their separate termination would not by
itself imply termination of their union.

\section{Conclusion}

The binary normal form isolates the contracting branch but does not itself
rule out strategic cycles. The missing global control is supplied by marked
finite outcome witnesses: strict token replacement gives an ordinal
projection, and the remaining equal-projection routing is well-founded by an
explicit normalizer analysis. Entry completeness then excludes every possible
draw continuation and proves finite optimal play from all odd starts.

\section*{Data, code, and correspondence}

The proof supplement, source code, tests, certificate files, Lean sources, and
reproducibility records are available at
\url{https://github.com/Grisha-Pochuev/3n-plus-minus-1-game}. The repository
is the authoritative location for the evolving verification artifacts; the
Zenodo record preserves the released article and source snapshot. Questions,
counterexamples, and audit reports may be sent to
\href{mailto:n_854@mail.ru}{n\_854@mail.ru}.

\bibliographystyle{plain}
\bibliography{references}

\end{document}
